\section{Experiments}
\label{sec:exp}
%Comment


%\subsection{Dataset}
%\label{subsec:dataset}
\noindent \textbf{Dataset.} 
We expand single-person datasets~\cite{zhang2021flow, xie2022vfhq, cui2025hallo3, zhu2022celebv, chung2018voxceleb2} with internet-collected data, yielding roughly 1,000 hours for first-stage training, and also gather two-person conversation clips for the second-stage training, retaining only about 12 hours after filtering.
Evaluations are conducted on two types of benchmarks: (i) standard talking-head benchmarks HDTF~\cite{zhang2021flow} and VFHQ~\cite{xie2022vfhq}, and (ii) our self-collected multi-person conversation dataset (head-and-body, both identities speak).  
We then select 20 videos from each benchmark, rigorously ensuring that their identities do not appear in the training set.
% Further details are provided in the supplementary material.

%\subsection{Experiment Settings}
\label{subsec:metric}

\noindent \textbf{Implement Details.}
To comprehensively evaluate our method, we train two models of different sizes: Wan2.1-1.3B-Inp and Wan2.1-I2V-14B~\cite{wan2025wan}, which serve as the foundational video diffusion models for our experiments. 
In all stages, the text~\cite{raffel2020exploring}, audio~\cite{baevski2020wav2vec}, and image~\cite{radford2021learning} encoders, as well as the 3D VAE, remain frozen.
The DiT main network, including the newly added AFCA layers, has all its parameters open for training. 
Stage 1 pretrains at $2 \times 10^{-5}$ learning rate; stage 2 fine-tunes at $5 \times 10^{-6}$.
All models are optimized with AdamW~\cite{loshchilov2017decoupled} on 32 NVIDIA H200 GPUs.
% The 14B model is trained using 32 NVIDIA H200 GPUs.
% In the first stage, the global batch size is set to 32, and the model is trained for 2.4M steps. 
% In the second stage, the batch size is adjusted to 16, and the model is trained for an additional 50K steps.
% The 1.3B model is trained using 8 NVIDIA H200 GPUs. 
% The global batch size is maintained at 48 throughout training, and the total number of training steps matches that of the 14B model.

\noindent \textbf{Evaluation Metrics.}
For the single-person benchmark, we employ several commonly used metrics: the Fréchet Inception Distance (FID)~\cite{heusel2017gans} and the Fréchet Video Distance (FVD)~\cite{unterthiner2019fvd} to assess the quality of the generated data, Sync-C~\cite{chung2016out} to measure the synchronization between audio and lip movements, and ID similarity~\cite{deng2019arcface} calculated between the first frame and the remaining frames.

For the multi-person benchmark, we evaluate from different dimensions. 
The newly introduced metric, termed \textit{Interactivity}, serves as the primary metric for assessment. 
For the FVD metric, the calculation is similar to that in the single-person benchmark. 
For the Sync-C metric, we refine its calculation as Sync-C$^{*}$ to focus only on the lip synchronization during each character's speaking periods, thereby avoiding the influence of long listening segments on the final lip synchronization score, specifically,
\begin{equation}
\text{Sync-C}^{*} = \frac{L_1 \cdot \text{Sync-C}_{L1} + L_4 \cdot \text{Sync-C}_{L4}}{L_1 + L_4}.
\end{equation}
Here, $L_1$ and $L_4$ denotes the speaking phases depicted in~\cref{fig:time}.

\noindent \textbf{Comparsion Methods.}
We compare AnyTalker with several state-of-the-art talking video generation methods.
For single-person generation, we compare with AniPortrait~\cite{wei2024aniportrait}, EchoMimic~\cite{chen2025echomimic}, Hallo3~\cite{cui2025hallo3}, Sonic~\cite{ji2025sonic}, FantasyTalking~\cite{wang2025fantasytalking}, StableAvatar~\cite{tu2025stableavatar}, OmniHuman-1.5~\cite{jiang2025omnihuman}, and MultiTalk~\cite{kong2025let}.
For multi-person generation, we choose Bind-Your-Avatar~\cite{huang2025bind} and MultiTalk~\cite{kong2025let} for quantitative and qualitative comparison.


\subsection{Comparison with SOTA methods}
% In this section, we comprehensively evaluate AnyTalker's ability to achieve a perfect balance among lip synchronization, video quality, and multi-person interactivity through both qualitative and quantitative experiments.
\noindent \textbf{Quantitative Comparison.}
To begin with, we compare AnyTalker with several single-person generation methods to verify its excellent single-person driving capability. 
The quantitative results are shown in~\cref{tab:single-benchmark}. 
Despite not being specifically designed for driving talking faces, AnyTalker achieves the best or competitive results across all metrics. 
Moreover, the 1.3B model of AnyTalker significantly outperforms AniPortrait~\cite{wei2024aniportrait}, EchoMimic~\cite{chen2025echomimic}, and StableAvatar~\cite{tu2025stableavatar} in terms of lip synchronization, even though they have a similar number of parameters. 
These results demonstrate the excellent and comprehensive driving capabilities of the AnyTalker framework.

Subsequently, we evaluate AnyTalker's ability to drive multiple IDs while maintaining both accurate lip synchronization and natural interactivity using the multi-person dataset, InteractiveEyes, described in~\cref{subsec:metric}, along with relevant metrics. 
In this comparison, we contrast AnyTalker with the available open-source multi-person driving methods, MultiTalk~\cite{kong2025let} and Bind-Your-Avatar~\cite{huang2025bind}. 
The results depicted in~\cref{tab:multi-benchmark} demonstrate that both the 1.3B and 14B models of AnyTalker achieve the best performance in terms of the \textbf{Interactivity} metric.
Additionally, the 14B model achieves the best results across all metrics, thereby validating the effectiveness of our proposed training pipeline.  
We further illustrate AnyTalker's capability to generate videos rich in interactivity through quantitative evaluation.


\begin{figure*}[t]
    \centering
    \includegraphics[width=\linewidth]{fig/qualitative.pdf}
    \caption{
    Qualitative comparison of multiple multi-person driving methods. 
    With the same text prompt, reference images, and multiple audio streams as input, we compare the generation results of Bind-Your-Avatar, MultiTalk, and AnyTalker. 
    The left case uses the input image from the InteractiveEyes dataset, while the right case uses the image produced by a text-to-image generative model~\cite{kolors}.
    }
    \label{fig:qualitative}
\end{figure*}

\noindent \textbf{Qualitative Comparsion.}
We then select an authentic human input from the InteractiveEyes dataset and use an input generated by an AIGC mode~\cite{kolors}, both accompanied by corresponding text prompts and dual audio streams, to conduct a quantitative evaluation comparison using Bind-Your-Avatar~\cite{huang2025bind}, MultiTalk~\cite{kong2025let}, and AnyTalker. 
As shown in~\cref{fig:qualitative}, AnyTalker generates more natural videos with eye and head interactions compared to the other methods. 
MultiTalk exhibits weaker eye interaction, while Bind-Your-Avatar tends to produce more static expressions. 
This trend further validates the effectiveness of the Interactivity metric proposed in~\cref{subsec:metric}. 
AnyTalker not only generates natural, two-person interactive speaking scenarios but also scales well to multiple IDs, as demonstrated in~\cref{fig:teaser}, where it effectively handles interactions among four IDs.
The qualitative results of the single-person benchmarks~\cite{zhang2021flow, xie2022vfhq} will be included in the supplementary material.


\begin{table}[t]
    \centering
    \caption{
        Ablation study about AnyTalker's components on the HDTF dataset using the 1.3B model.
        ``Baseline'' denotes the model equipped solely with the basic audio attention layers. 
        ``AFCA'' indicates the inclusion of the Audio-Face Cross Attention mechanism. 
        ``Single'' indicates that authentic multi-person data is not utilized in this stage.
    }
    \small
    \begin{tabular}{cccccc}
        \toprule
        \multirow{2}{*}{\textbf{Setting}} & \multicolumn{4}{c}{\textbf{Metrics}} \\
        \cmidrule(lr){2-5}
        & \textbf{Sync-C$\uparrow$} & \textbf{FID} & \textbf{FVD} & \textbf{ID$\uparrow$} \\
        \midrule
        Baseline & 5.42 & \textbf{13.01} & \underline{170.96} & \underline{0.90} \\
        w/o AFCA & \underline{6.71} & 14.97 & 207.47 & 0.88 \\
        w/o Mask Token & 5.84 & 14.81 & 193.78 & 0.89 \\
        w/o Concatenated Data & 6.21 & \underline{14.73} & 202.01 & \textbf{0.91} \\
        AnyTalker-1.3B (Single) & \textbf{6.97} & 15.58 & \textbf{166.27} & \textbf{0.91} \\
        \bottomrule
    \end{tabular}
    \label{tab:ab-component}
\end{table}

\subsection{Ablation Study}
\label{subsec:ab}

\noindent \textbf{Components.}
We conduct ablation studies on the three important components mentioned in~\cref{fig:method}, including Audio-Face Cross Attention, Mask Token for attention output, and concatenated multi-person data.
Starting from the full 1.3B model in the first stage, which utilizes only single-person data, we progressively remove the relevant components to evaluate their impact on lip synchronization, generation quality, and identity similarity in the generated videos. 
% The results are shown in~\cref{tab:ab-component}. 
% Each component plays a significant role in each evaluation dimension, especially the Concatenated Multi-Person Data, which greatly enhances lip movement accuracy. 
The results in~\cref{tab:ab-component} show that every component plays a significant role, with concatenated multi-person data most beneficial to lip-movement accuracy.
Although the completed model in the initial phase exhibits a marginally higher FID score compared to the Baseline model, we attribute this to the Baseline model's propensity for generating talking videos with minimal facial expressions or head movements. 
Conversely, the completed model generates more dynamic actions, which inherently influence the FID calculation to some degree. 
We regard this trade-off as a natural consequence of the model's design. Furthermore, the completed model demonstrates superior performance over the Baseline model across other evaluation metrics. 
Crucially, the first-stage model already has the basic capability to drive multiple individuals, a feature that the Baseline model lacks.

\begin{table}[t]
    \centering
    \small
    \caption{
        Ablation studies conducted on the InteractiveEyes dataset using the 1.3B model. 
        ``RS'' denotes the use of authentic single-person data in the first stage. 
        ``CM'' indicates concatenated multi-person data in the first stage. 
        ``RM'' represents authentic multi-person data in the second stage.
    }
    \begin{tabular}{cccccc}
        \toprule
        \multicolumn{3}{c}{\textbf{Setting}} & \multicolumn{3}{c}{\textbf{Metrics}} \\
        \cmidrule(lr){1-3} \cmidrule(lr){4-6}
        \textbf{RS} & \textbf{CM} & \textbf{RM} & \textbf{Interactivity$\uparrow$}& \textbf{Sync-C$^{*}$$\uparrow$} & \textbf{FVD}  \\
        \midrule
        $\times$ & $\checkmark$ & $\times$ & 0.55 & 3.21 & 672.18  \\
        $\checkmark$ & $\times$ & $\times$ & 0.47 & 4.13 & 475.31  \\
        $\checkmark$ & $\checkmark$ & $\times$ & 0.58 & \textbf{4.89} & \textbf{393.86}  \\
        \midrule
        $\checkmark$ & $\times$ & $\checkmark$ & \underline{0.71} & 3.63 & 511.51 \\
        $\checkmark$ & $\checkmark$ & $\checkmark$ & \textbf{0.97} & \underline{4.56} & \underline{467.84}  \\
        \bottomrule
    \end{tabular}
    \label{tab:ab-data}
\end{table}



\noindent \textbf{Multi-Person Data.}
Subsequently, we focus on multi-person data and conduct more targeted ablation experiments on the InteractiveEyes benchmark. 
In the first stage, we create concatenated multi-person data using single-person data. As shown in~\cref{tab:ab-data}, models utilizing this concatenated data outperform those without it across all evaluation dimensions, especially in lip synchronization and interactivity, highlighting the role of concatenated data in learning multi-person speaking patterns. 
It is worth noting that the results in the first row of the table also demonstrate the necessity of mixing single-person data in the first stage; otherwise, the generated results would be unstable.

After fine-tuning with authentic multi-person data, models that use concatenated data in the first stage demonstrate even more superior performance, further proving their early adaptation to multi-person speaking patterns. 
Although fine-tuning with authentic multi-person data leads to a slight decrease in lip synchronization, it significantly improves interactivity, which we consider a reasonable trade-off.

